\scanpage{149}
\deleted{ou plus simplement, si $\theta=\exp2i\pi b$, en termes de $a$ et $b$
\[
(36)\qquad\left\{\begin{array}{rcll}
&\smash{\raisebox{-1ex}{$\Big\{$}}&a\equiv\tfrac12,&b\equiv1\quad(\mathbb Z)\\
\text{ou}&&a\equiv1,&b\equiv\tfrac12\quad(\mathbb Z)
\end{array}\right.
\]
Nous nous intéressons plus particulièrement au sous-groupe}

Nous nous intéressons au sous-groupe $G_0$ de $G$ correspondant à la valeur $\theta=1$, on écrira $g_a$ au lieu de $g_{a,\theta}$, donc
\[
(34)\qquad g_a(\unsure{z},s)=(z\exp2i\pi\unsure{as},s)
\]
Pour que cette opération induise l'identité sur $\partial X$, il f. et s. qu'on ait
\[
(35)\qquad\exp2i\pi a=1\qquad\text{i.e.}\quad a\in\mathbb Z.
\]
Nous nous intéressons au sous-groupe
\margin{149}
